\chapter{Структура проекта и артефактов}
\label{appendix:artifacts}

В приложении приведен пример структуры файлов, используемой в MVP-пайплайне. Она отражает путь от исходных документов до векторного индекса и отчетов.

\begin{lstlisting}[language=bash,caption={Основные каталоги и артефакты MVP}]
data/
  raw/                         # исходные DOCX, PDF, XLSX
  registry/
    documents_registry.csv     # реестр документов
  parsed/
    documents.jsonl            # parsed-корпус
  processed/
    chunks.jsonl               # чанки
    chunk_embeddings.npy       # матрица embeddings
    chunk_embeddings_manifest.jsonl
  reports/
    parsing/parse_report.json
    parsing/parsed_corpus_validation.json
    chunking/chunking_report.json
    embeddings/embedding_report.json
    indexing/qdrant_index_report.json
    retrieval/retrieval_eval_report.json
  eval/
    retrieval_questions.jsonl
    rag_questions.jsonl

scripts/
  generate_documents_registry.py
  parse_documents.py
  validate_parsed_corpus.py
  chunk_documents.py
  embed_chunks.py
  index_qdrant.py
  search_qdrant.py
  answer_question.py

src/robymaxrag/
  parsing/
  chunking/
  embeddings/
  retrieval/
  generation/
\end{lstlisting}

Такая структура позволяет воспроизводить каждый этап отдельно и использовать отчеты для контроля качества преобразования данных.

%%% Local Variables:
%%% TeX-engine: xetex
%%% End:
